\section{The sharp constant of the RMS lower bound}
\label{app:sharp-rms}

The proof of \cref{thm:main} in \cref{sec:minimax-proof} gives the lower
constant $1/64$. The best constant is $1/6$.

\begin{theorem}\label{thm:sharp-rms}
For every probability law $\mu$ on $(0,1/2]$ and every $B\ge1$,
\begin{equation}\label{eq:sharp-rms}
 R_\mu(B)\ \ge\ \tfrac16\,\om(B^{-1/2}).
\end{equation}
The constant cannot be increased, since $R_{\delta_{1/2}}(4)=1/6$. For
$B\ge25$ it improves to $\sqrt{59/2100}>0.1676$. The same three statements
hold if every coin's magnitude is revealed for free, and if the algorithm
selects the magnitude of each fresh coin.
\end{theorem}

The experiment is that of \cref{thm:main}. An algorithm may use private
randomness, interleave and revisit coins, and stop at an unbounded time.
The budget bounds its expected work under every sign law.

\paragraph{Formal scope.}
\sharprmsleanref{Full.lean}{sharp\_lower} proves \eqref{eq:sharp-rms}
against the same experiment, risk and modulus definitions as the formal
statement of \cref{thm:main}.
\sharprmsleanref{Full.lean}{sharp\_coefficient} proves that a constant is
valid for every law and every $B\ge1$ exactly when it is at most $1/6$,
and \sharprmsleanref{Full.lean}{large\_budget\_lower} gives the constant
for $B\ge25$. The oracle forms are
\sharprmsleanref{Full.lean}{oracle\_sharp\_lower},
\sharprmsleanref{Full.lean}{oracle\_large\_budget\_lower\_bound} and
\sharprmsleanref{Full.lean}{oracle\_endpoint}. The Lean proof follows the
route below. Among its differences: for $B\le4$ it integrates the exact
Bayes risk under the uniform prior; it proves the bound on $J(t)$ by an
inequality at each stopped transcript, summed over the coins; it proves
\cref{lem:sharp-scale} by algebra on each exit; and it proves
\eqref{eq:sharp-scalar} by a polynomial bound. The Lean kernel evaluates
the integer comparisons behind \eqref{eq:sharp-window}.
\Cref{thm:sharp-bracket} in \cref{sec:sharp-bracket} narrows the best
constant for $B\ge25$ to the interval $[0.175198,0.186752]$.

Throughout, a flip is $Y\in\{-1,1\}$, and a coin of gap $d$ and sign
$S\in\{-1,1\}$ has $\E[Y\mid S]=2dS$. Write $a=2d$ for its bias. A sign
law is described by its sign mean $\E[S\mid d]=2q(d)-1$.

\subsection{Reduction to one gap}

\begin{lemma}\label{lem:sharp-reduction}
For every gap $g\in(0,1/2]$ and every $B>0$,
\[
 R_\mu(B)\ \ge\ \om(g)\,R_{\delta_g}(B).
\]
The same holds with either oracle risk on the left.
\end{lemma}

\begin{proof}
Put $s=2g$ and $f(d)=\min\{1,s/a\}$, so that $\om(g)=\int f\,d\mu$.
Consider only the sign laws with sign mean $tf(d)$, $-1\le t\le1$. Their
target is $(1+t\,\om(g))/2$. Call a coin of gap $g$ whose sign $S$ has mean
$t$ a source coin.

Run a given protocol for $\mu$ on source coins. Use one fresh source coin
for each coin the protocol requests and one source flip for each flip it
requests. Draw the magnitude of each requested coin as the protocol's
experiment prescribes: from $\mu$, or from the revealed array, or by the
protocol's own selection. It remains to turn source flips into flips of a
coin of bias $a$ whose sign has mean $tf$.

If $a\le s$, then $f=1$. Keep each source flip with probability
$(1+a/s)/2$ and negate it otherwise, independently. Given $S$ the results
are independent with mean $aS$.

If $a>s$, then $f=s/a$. The required coin has a sign $S'$ with
$\E[S'\mid S]=fS$, which the simulation never draws. Suppose the simulated
flips so far have net count $K$, the sum of their values, and put
$z=\tanh(K\artanh a)$ if $a<1$.
Given $S$ and these flips, the next flip of the required coin has mean
\[
 a\,\frac{z+fS}{1+fzS}=b+c\,sS,\qquad
 b=\frac{az(1-f^2)}{1-f^2z^2},\qquad c=\frac{1-z^2}{1-f^2z^2}.
\]
If $z=0$, return the source flip. Otherwise return it with probability
$c$, and with probability $1-c$ return an independent sign of mean $a/z$.
This is possible because $|z|\ge a$ whenever $K\ne0$, and it has the
required mean because $(1-c)a/z=b$. If $a=1$, the first source flip has
mean $sS=fS$ given $S$: use it as $S'$ and repeat it at every revisit.

A flip is determined in law by its conditional mean, and each step uses a
fresh source flip and fresh randomness. Hence the simulated transcript has
the law of the original experiment at the sign law with sign mean $tf$,
under any interleaving, and the same work. If $\widehat\theta$ is the
protocol's output, then $\tfrac12+(2\widehat\theta-1)/(2\,\om(g))$
estimates the source target $(1+t)/2$ with the RMS error of
$\widehat\theta$ divided by $\om(g)$. The simulation is a protocol for
$\delta_g$ with budget $B$.
\end{proof}

The Lean proof
\sharprmsleanref{Reduction.lean}{critical\_gap\_reduction} states the lemma
for the minimax risks, and
\sharprmsleanref{OracleReduction.lean}{oracle\_gap\_reduction} for an
arbitrary measurable array of magnitudes given the oracle seed.

For $B\ge4$ take $g=1/\sqrt B$. Then $\om(g)=\om(B^{-1/2})$, and the gap
has bias $s=2/\sqrt B\le1$ with $s^2B=4$. Write
\[
 r_s=R_{\delta_{s/2}}(4/s^2),\qquad 0<s\le1.
\]
\Cref{thm:sharp-rms} for $B\ge4$ follows from \cref{lem:sharp-reduction}
and the next proposition. \Cref{sec:sharp-small} treats $B\le4$.

\begin{proposition}\label{prop:sharp-critical}
For every $0<s\le1$, $r_s\ge1/6$. If $s\le2/5$, then
$r_s\ge\sqrt{59/2100}$.
\end{proposition}

\subsection{The single-gap experiment}\label{sec:sharp-single}

Fix $0<s<1$ and $B=4/s^2$. Every coin has bias $s$, the sign mean
$t\in[-1,1]$ is unknown, and the target is $(1+t)/2$. An estimator $X$ of
$t$ has four times the squared risk of $(1+X)/2$. Let $P_t$ be the law of
the transcript, including the seed.

Clip $X$ to $[-1,1]$, stop after $M$ flips and return a fixed value on
overflow. The squared risk for $t$ increases by at most $4B/M$, uniformly
in $t$. All calculations below are therefore made at a bounded total
horizon.

At a stopped transcript let $K_i$ be the net count of coin $i$ and
$z_i=\tanh(K_i\artanh s)$. Under $P_0$ the signs are fair and $z_i$ is the
conditional mean of $S_i$. For every $t$,
\begin{equation}\label{eq:sharp-likelihood}
 \frac{dP_t}{dP_0}=\prod_i(1+tz_i).
\end{equation}
Given the signs and the past, a paid flip $Y$ of coin $i$ has
$\E[S_iY]=s$. Summing over paid flips and conditioning on the transcript
under $P_0$ gives
\begin{equation}\label{eq:sharp-wald}
 \E_0\sum_i z_iK_i=s\,\E_0T .
\end{equation}
Both identities hold when a revisit depends on other coins.

\subsection{Large budgets: three sign biases}\label{sec:sharp-large}

Let $s\le2/5$ and $\beta=7/10$. By \eqref{eq:sharp-likelihood},
\[
 \KL(P_0\Vert P_\beta)+\KL(P_0\Vert P_{-\beta})
 =\E_0\sum_i-\log(1-\beta^2z_i^2).
\]
For $|z|<1$, $-\log(1-\beta^2z^2)\le\beta^2z\artanh z$: the coefficients
of $z^{2j}$ are $\beta^{2j}/j$ and $\beta^2/(2j-1)$, and
$\beta^{2j-2}\le1/2<j/(2j-1)$ for $j\ge2$. Since
$z_i\artanh z_i=z_iK_i\artanh s$, \eqref{eq:sharp-wald} and
$\E_0T\le4/s^2$ give
\begin{equation}\label{eq:sharp-cost}
 \KL(P_0\Vert P_\beta)+\KL(P_0\Vert P_{-\beta})
 \le\beta^2s\artanh(s)\,\E_0T
 \le4\beta^2\,\frac{\artanh s}{s}
 \le\frac{3283}{1575}.
\end{equation}
The last step uses
$\artanh(s)/s\le1+s^2/(3(1-s^2))\le67/63$ for $s\le2/5$.

For the risk, put
$\phi_\pm(x)=1+8\bigl((x\mp\beta)^2-1/9\bigr)\ge1/9$. For probability
laws $P,Q$ and a positive function $\phi$,
$\E_P\log\phi\le\KL(P\Vert Q)+\E_Q\phi-1$. With $P=P_0$,
$Q=P_{\pm\beta}$ and $\phi=\phi_\pm(X)$ this gives
\[
 \KL(P_0\Vert P_{\pm\beta})\ge\E_0\log\phi_\pm(X)
   -8\,\E_{\pm\beta}(X\mp\beta)^2+\tfrac89 .
\]
The product is $\phi_+(X)\phi_-(X)=G(X^2)$ with
$G(y)=(907/225+8y)^2-(3136/25)\,y$, and for $y\ge0$
\begin{equation}\label{eq:sharp-scalar}
 \tfrac12\log G(y)+\tfrac52\,y\ \ge\ \tfrac43 .
\end{equation}
Adding the two bounds and using \eqref{eq:sharp-scalar} and
\eqref{eq:sharp-cost},
\[
 \tfrac{40}9\le5\,\E_0X^2+8\,\E_\beta(X-\beta)^2+8\,\E_{-\beta}(X+\beta)^2
 +\tfrac{3283}{1575}.
\]
If the three squared risks are at most $\varrho^2$, then
$21\varrho^2\ge40/9-3283/1575$, that is, $\varrho^2\ge59/525$. Hence
$r_s^2\ge59/2100>1/36$.

To prove \eqref{eq:sharp-scalar}, note that its derivative has the sign of
$G'(y)+5G(y)=320y^2-\tfrac{7952}{45}y+\tfrac{205609}{10125}$, which is
positive at $0$, negative at $19/50$ and positive at $39/100$. The minimum
is at $0$ or at the larger root, which lies in $[19/50,39/100]$. At $0$
the value is $\log(907/225)>\log4>4/3$. On that interval $G$ decreases
and $G(39/100)=(67/45)^2$, so the value at the larger root is at least
$\log(67/45)+19/20\ge11/28+19/20>4/3$, by $\log x\ge2(x-1)/(x+1)$ for $x\ge1$.
\sharprmsleanref{WeakPopulation.lean}{weak\_gap\_minimax\_lower} is the
Lean statement of this case.

\subsection{Intermediate budgets: the exact information price}
\label{sec:sharp-middle}

Let $J(t)$ be the Fisher information of the transcript and $v=1-t^2$.
Give $t$ the uniform prior on $(-1,1)$ and write $\overline\E$ for the
average over $t$ and the transcript, and $\xi_t$ for the score. For a
bounded-horizon estimator $X$ with values in $[-1,1]$, integration by
parts gives
\begin{equation}\label{eq:sharp-score}
 \overline\E\bigl[(X-t)(v\xi_t-2t)\bigr]=\tfrac23,\qquad
 \overline\E\bigl[(v\xi_t-2t)^2\bigr]
 =\tfrac43+\tfrac12\int_{-1}^1v^2J(t)\,dt .
\end{equation}
The boundary term vanishes because $v(\pm1)=0$. If the second quantity is
at most $4$, the Cauchy--Schwarz inequality gives
$\overline\E(X-t)^2\ge1/9$, hence $r_s\ge1/6$.

\paragraph{One coin.}
Fix $0\le t<1$ and a stopping rule for one coin, with net count $K$ and
$Z=\tanh(K\artanh s)$ at its stop. Let $W_t$ be its expected number of
flips and $I_t$ its Fisher information about $t$. By
\eqref{eq:sharp-likelihood} and the argument for \eqref{eq:sharp-wald},
\begin{equation}\label{eq:sharp-local}
 \E_0Z=0,\qquad I_t=\E_0\frac{Z^2}{1+tZ},\qquad
 sW_t=\E_0\bigl[(Z+t)K\bigr].
\end{equation}
These identities are proved at a bounded horizon. They extend to every
rule with finite expected work, by stopping it after $n$ flips and letting
$n\to\infty$.

Put $x_k=\tanh(k\artanh s)$. The rule that stops when the net count
reaches $+1$ or $-k$ ends at $Z=s$ or $Z=-x_k$, with fair probabilities
$x_k/(s+x_k)$ and $s/(s+x_k)$. Its work has finite mean, since $k+1$ equal
flips in a row end it. Substituting in \eqref{eq:sharp-local},
\begin{equation}\label{eq:sharp-exit}
 \eta_k(s,t)=\frac{I_t}{s^2W_t}
 =\frac{1+x_k/s}{(1-tx_k)(1+ts)\bigl[k+1-t(k/x_k-1/s)\bigr]},
 \qquad \eta_s(t)=\sup_{k\ge1}\eta_k(s,t).
\end{equation}
Revealing the sign bounds $I_t$ by $1/v$, and $W_t\ge1$, so
$\eta_s(t)\le1/(s^2v)$.

\begin{lemma}\label{lem:sharp-price}
Every stopping rule for one coin satisfies $I_t\le s^2\eta_s(|t|)\,W_t$.
\end{lemma}

\begin{proof}
Let $t\ge0$ and $\lambda=1/\eta_s(t)$. The exit with $k=1$ gives
$\lambda\le1-s^2t^2\le1$.
For a net count $K\ne0$ with $z=\tanh(K\artanh s)$ put
\[
 F(K)=\frac{s(z+t)K}{z}-\frac{\lambda z}{1+tz}.
\]
By \eqref{eq:sharp-local}, $s^2W_t-\lambda I_t=\E_0[ZF(K)]$, with the term at
$K=0$ equal to zero, and $\E_0[ZF(K)]=\E_0[Z(F(K)-F(1))]$ because
$\E_0Z=0$. For $K\ge1$ the first term of $F$ is $sK+stK/z$. It increases
by at least $s$ at each step, since $K/z$ increases. The second term
increases by at most $s$, since $z$ increases by at most $s$ and
$z\mapsto\lambda z/(1+tz)$ has slope at most one. Thus $F(K)\ge F(1)$. For
$K=-k$, the exit at $+1$ or $-k$ satisfies $\lambda I_t\le s^2W_t$ by the
definition of $\lambda$, which reads $F(-k)\le F(1)$. Hence
$Z(F(K)-F(1))\ge0$
at every reachable count. A negative $t$ follows by reversing the flips,
and a randomized rule by averaging.
\end{proof}

\paragraph{Any adaptive protocol.}
By \eqref{eq:sharp-likelihood}, $J(t)=\E_t\sum_iz_i^2/(1+tz_i)^2$. For each
coin, condition on the seed and on the transcript before its first flip.
Each other coin has a sequence of outcomes that its later flips would
give, whether or not the protocol requests them. Given that past and at
the same $t$, these sequences are independent of the coin's sign and
flips. With them as private randomness, the protocol is a randomized rule
for the one coin. Its information is the expectation of the coin's
summand and its work is the number of flips paid for the coin.
\Cref{lem:sharp-price} and the budget give
\[
 J(t)\le s^2\eta_s(|t|)\,\E_tT\le4\eta_s(|t|).
\]
This is \sharprmsleanref{AdaptivePrice.lean}{adaptive\_fisher\_price}. By
\eqref{eq:sharp-score} it suffices to prove that, for $2/5\le s<1$,
\begin{equation}\label{eq:sharp-window}
 \int_0^1(1-t^2)^2\eta_s(t)\,dt\ \le\ \tfrac23 .
\end{equation}

\begin{lemma}\label{lem:sharp-scale}
If $0<s\le u<1$, then $s^2\eta_s(t)\le u^2\eta_u(t)$.
\end{lemma}

\begin{proof}
Negate each flip of a coin of bias $u$ independently with probability
$(1-s/u)/2$. Given the sign, the results are the flips of a coin of bias
$s$. A stopping rule for the latter becomes a randomized rule for the
former with the same work, and keeping the original flips can only
increase information. Apply \cref{lem:sharp-price} at bias $u$.
\end{proof}

\paragraph{Near-perfect coins.}
Let $39/40\le s<1$. The exit with $k=1$ has $\eta_1(s,t)=1/(1-s^2t^2)$. For
$k\ge2$ the rule pays for a second flip when the first is negative, so
$W_t\ge1+(1-st)/2$ and $\eta_k(s,t)\le1/\bigl(s^2v[1+(1-st)/2]\bigr)$. For
$t\le3/4$,
\[
 s^2v\,\frac{1-st}2\ \ge\ \Bigl(\frac{39}{40}\Bigr)^2\frac7{16}\cdot\frac18
 \ >\ 1-\Bigl(\frac{39}{40}\Bigr)^2\ \ge\ 1-s^2,
\]
which gives $\eta_k(s,t)\le \eta_1(s,t)$, so $\eta_s(t)=\eta_1(s,t)$ there. For
$t>3/4$ use $\eta_s(t)\le1/(s^2v)$. Compared with $v$, the integrand
$v^2\eta_1(s,t)$ is smaller by
at least $(1-s^2)vt^2$ on $[0,3/4]$, and $v/s^2$ is larger by
$(1-s^2)v/s^2$ on $[3/4,1]$. Since
\[
 \int_0^{3/4}vt^2\,dt=\frac{477}{5120}
 >\Bigl(\frac{40}{39}\Bigr)^2\frac{11}{192}
 \ge s^{-2}\int_{3/4}^1v\,dt
\]
and $\int_0^1v\,dt=2/3$, \eqref{eq:sharp-window} holds.

\paragraph{The window $2/5\le s\le39/40$.}
Here \eqref{eq:sharp-window} is checked by integer arithmetic. For a
rational bias, $x_1=s$ and $x_{k+1}=(x_k+s)/(1+sx_k)$ are rational. Each
$\eta_k(s,t)$ increases with $t\ge0$. The factor $k+1-t(k/x_k-1/s)$ is
positive and increases with $k$. For $k\ge20$, replace it by its value at
$k=20$ and replace $x_k$ by $1$ elsewhere in \eqref{eq:sharp-exit}; the
result bounds $\eta_k(s,t)$ and also increases with $t$. On each of $2048$
equal subintervals of $[0,1]$ let $v_0$ be the value of $v$ at the left
endpoint. Multiply each of these twenty expressions at the right endpoint
by $v_0^2$, round up to a multiple of $2^{-32}$ and take the largest. The
height of the subinterval is the smaller of this and $v_0/s^2$, also
rounded up. On the last subinterval the twentieth expression is infinite
at the right endpoint, and the height is $v_0/s^2$. Let $U(u)$ be the mean of these heights for the bias $u$; it
bounds the integral in \eqref{eq:sharp-window} at $s=u$. By
\cref{lem:sharp-scale}, $(u/\ell)^2U(u)\le2/3$ proves
\eqref{eq:sharp-window} for every $s\in[\ell,u]$. This inequality holds on $140$ adjacent intervals that
cover $[2/5,39/40]$, with endpoints that are multiples of $1/40960$; the
largest value is below $0.66666$. The comparisons are
\sharprmsleanref{Window.lean}{checked\_window}, and
\sharprmsleanref{MiddleBudget.lean}{critical\_minimax\_lower} is the
resulting bound $r_s\ge1/6$ for every $0<s\le1$.

\subsection{Budgets at most four, and equality}\label{sec:sharp-small}

Let $B\le4$ and let $\mu$ be arbitrary. Consider the sign laws with the
same sign mean $t$ at every magnitude; the target is $(1+t)/2$. Reveal the
sign of a coin at its first paid flip. This can only help the estimator.
A fresh sign has information $1/v$ about $t$, and a revisit has none, so
$J(t)\le B/v$ at a bounded horizon. The identities
\eqref{eq:sharp-score} hold in this experiment by the same integration by
parts, and
$\tfrac43+\tfrac12\int_{-1}^1v^2J\le\tfrac43+\tfrac23B\le4$. Hence
$R_\mu(B)\ge1/6\ge\om(B^{-1/2})/6$ for every $B\le4$, and $r_1\ge1/6$.
This completes the proofs of \cref{prop:sharp-critical} and of
\eqref{eq:sharp-rms}. The Lean statement is
\sharprmsleanref{FourBudget.lean}{four\_budget\_minimax\_lower}.

For equality take $\mu=\delta_{1/2}$ and $B=4$, so that
$\om(B^{-1/2})=1$. One flip reveals the sign of a fresh coin. If $N$ is
the number of positive signs among four fresh coins and $q$ is the
fraction of positive coins, then $(N+1)/6$ has squared risk
\[
 \frac{4q(1-q)+(1-2q)^2}{36}=\frac1{36}
\]
for every $q$. Hence $R_{\delta_{1/2}}(4)=1/6$, which is
\sharprmsleanref{FourBudget.lean}{endpoint\_minimax\_eq}. The same four
flips are available to both oracles. Their lower bounds follow from
\cref{lem:sharp-reduction} and the bounds above for one gap: take
$g=1/\sqrt B$ if $B\ge4$, and $g=1/2$ if $B\le4$, where $\om(1/2)=1$.

\begin{remark}
Equality in \eqref{eq:sharp-rms} requires $B=4$: every case of the proof
is strict for $B\ne4$. The upper constant of \cref{thm:main} is not
improved, so the two sides of \eqref{eq:main-rms} remain a factor $3072$
apart.
\end{remark}

\subsection{Budgets of at least 25}\label{sec:sharp-bracket}

\Cref{sec:sharp-large} gives the constant $\sqrt{59/2100}>0.1676$ for
$B\ge25$. Two further arguments narrow the best constant for these
budgets. Lean checks both. The text below outlines the Lean proofs; their
numerical certificates are checked there in exact rational and integer
arithmetic and are not repeated here.

\begin{theorem}\label{thm:sharp-bracket}
Let $c_{25}$ be the largest constant such that
$R_\mu(B)\ge c_{25}\,\om(B^{-1/2})$ for every law $\mu$ and every
$B\ge25$. Then
\[
 0.175198\ \le\ c_{25}\ \le\ 0.186752 .
\]
The lower bound also holds for both oracle risks. For $B\ge100$ the
constant $0.17807955$ holds in all three models.
\end{theorem}

\paragraph{Formal scope.}
\largermsleanref{Full.lean}{bracket\_lower} proves the lower bound in the
three models, \largermsleanref{Full.lean}{hundred\_lower} the bound for
$B\ge100$, \largermsleanref{Full.lean}{fifth\_gap\_upper} the bound
$R_{\delta_{1/5}}(25)\le0.186752$, and
\largermsleanref{Full.lean}{bracket\_upper} the upper bound on $c_{25}$.
They use the same definitions as \cref{thm:sharp-rms}.

\paragraph{The lower bound.}
By \cref{lem:sharp-reduction} with $g=1/\sqrt B$ it suffices to bound
$r_s$ for $s=2/\sqrt B\le2/5$, in the experiment of
\cref{sec:sharp-single} and at a bounded horizon. For $t^2\le2/3$,
\begin{equation}\label{eq:sharp-paired}
 \KL(P_0\Vert P_t)+\KL(P_0\Vert P_{-t})
 \le-\log(1-t^2s^2)\,\E_0T
 \le-B\log(1-4t^2/B).
\end{equation}
Lean proves the first inequality from \eqref{eq:sharp-likelihood} and
\eqref{eq:sharp-wald}, by bounding each coin's term through its
posterior. It replaces \eqref{eq:sharp-cost}.

Take five biases $t_i$, namely $0.527436$, $0.679980$, $0.753283$,
$0.798023$ and $0.816490$, slopes $b_i$, namely $6.121167$, $7.070320$,
$7.599955$, $7.966557$ and $8.119407$, and put $a=19.780099$ and
$R=4\cdot0.175198^2$. Write
\[
 \pi(B)=-B\sum_{i=1}^5\log\bigl(1-4t_i^2/B\bigr),\qquad
 D=a+2\sum_{i=1}^5b_i=93.534911 .
\]
Clip the estimate $X$ of the sign mean to $[-1,1]$. The tests
$g_{i,\pm}(x)=1+b_i\bigl((x\mp t_i)^2-R\bigr)$ are positive. Apply
$\E_P\log\phi\le\KL(P\Vert Q)+\E_Q\phi-1$ with $P=P_0$,
$Q=P_{\pm t_i}$ and $\phi=g_{i,\pm}(X)$, sum over the ten alternatives
and use \eqref{eq:sharp-paired}:
\[
 \E_0\Bigl[aX^2+\sum_{i,\pm}\log g_{i,\pm}(X)\Bigr]
 \le a\,\E_0X^2+\pi(B)
 +\sum_{i,\pm}b_i\bigl(\E_{\pm t_i}(X\mp t_i)^2-R\bigr).
\]
Lean checks the pointwise inequality
\begin{equation}\label{eq:sharp-five}
 a x^2+\sum_{i,\pm}\log g_{i,\pm}(x)\ \ge\ aR+\pi(25),\qquad |x|\le1,
\end{equation}
with a Taylor lower bound of degree $27$ for the exponential and exact
polynomial checks on nine intervals in $x^2$. If the eleven squared
risks $\E_0X^2$ and $\E_{\pm t_i}(X\mp t_i)^2$ are at most $S$, the two
displays give $D(S-R)\ge\pi(25)-\pi(B)$. The risk for the target is half
the RMS error of $X$, so
\[
 r_s\ \ge\ \frac12\sqrt{R+\frac{\pi(25)-\pi(B)}{D}} .
\]
The function $\pi$ decreases in $B$. The right side is therefore at
least $\sqrt R/2=0.175198$ for $B\ge25$, and at least $0.17807955$ for
$B\ge100$. \Cref{lem:sharp-reduction} carries both bounds to every law
and to both oracles.

\paragraph{The upper bound.}
Take $\mu=\delta_{1/5}$ and $B=25$, so that $s=2/5$ and
$\om(B^{-1/2})=1$. Flip $18$ fresh coins and let $j$ be the number of
flips equal to $+1$. Then make seven more flips. Before the $(k+1)$st of
them, a fixed table indexed by $j$ and the $k$ earlier answers chooses a
fresh coin, the $(k+1)$st coin whose first flip was $+1$, or the
$(k+1)$st coin whose first flip was $-1$, counted in the order of the
first $18$ flips. A missing coin is replaced by a fresh one. The estimate
is a rational with denominator $100{,}000$, indexed by $j$ and the seven
answers. Every path uses at most $25$ flips. The table was found by
backward minimization of the Bayes squared loss under a symmetric
four-point prior; only the risk of the resulting policy is checked.

If a fraction $p$ of the coins have positive sign, a fresh flip is $+1$
with probability $(3+4p)/10$. Two flips of one coin are both $+1$ with
probability $(9+40p)/100$, and they differ with probability $21/100$ in
either order. The squared risk of the estimate for the target $p$ is
therefore a polynomial in $p$. Lean computes it exactly and compares it
with its value at the endpoints on eight rational subintervals of
$[0,1]$. Its maximum,
\[
 348763093753675683898541718818533\cdot10^{-34}\ <\ 0.186752^2,
\]
is attained at $p=0$ and at $p=1$. Hence $R_{\delta_{1/5}}(25)\le0.186752$,
and no constant above $0.186752$ is valid at $\mu=\delta_{1/5}$ and
$B=25$. The value of $c_{25}$ remains open.
